\SourceLeaf{306}
\Needspace{5\baselineskip}\begin{center}
[2.]
\end{center}

If $B, C$ are two points on $a$, one may quite properly set $B - C = ai$.

\Needspace{5\baselineskip}\begin{center}
[3.]
\end{center}

The main question now is what forces should be applied to the vertices of the pentagon in the barycentric calculus or to its sides in the calculus of resultants, so that the relations between those forces and their equivalents on the inner pentagon are as simple as possible.

\Needspace{5\baselineskip}\begin{center}
[4.]
\end{center}

The collinearity of two systems of points or straight lines is perhaps best defined by the fact that one and the same linear equation can always be set up between any 5 points (lines) of one system and the 5 corresponding ones of the other.

\Needspace{5\baselineskip}\begin{center}
[5.]
\end{center}

The characteristic feature of a collineation seems to be most simply expressed as follows.

Let $A, B.\,.\,.$ be straight lines and $P,Q.\,.\,.$ points, all in one system, corresponding in the other to the lines $a,b.\,.\,.$ and the points $p,q.\,.\,.$. If $PA$ denotes the (perpendicular) distance of the point $P$ from the straight line $A$, then in general
\[
\frac{PA.QB.pb.qa}{pa.qb.PB.QA} = 1.
\]
